% Modeled per-transaction cost as the control-plane rule set grows.
% Expressive engines branch on business cases on the hot path (cost rises
% with rule-set size); the conservation kernel does a fixed amount of
% integer work per fixed-width record (flat). Analytical model, not a
% measured benchmark --- see Section "Analysis and Value Proposition".
\begin{figure}[t]
\centering
\begin{tikzpicture}
\begin{axis}[
    width=\linewidth, height=6cm,
    xmode=log, log basis x=2,
    xtick={1,4,16,64,256,1024},
    xticklabels={1,4,16,64,256,1024},
    xlabel={control-plane rule-set size (number of rules)},
    ylabel={modeled per-tx cost (ns)},
    ymin=0,
    legend style={at={(0.02,0.97)}, anchor=north west, draw=none,
                  fill=white, fill opacity=0.85, text opacity=1, font=\scriptsize},
    grid=both, grid style={black!10}, mark options={solid},
]
\addplot[mark=*, thick, red!60!black]
    table[x=rules, y=expressive] {figures/data/throughput.dat};
\addplot[mark=square*, thick, green!45!black]
    table[x=rules, y=kernel] {figures/data/throughput.dat};
\legend{expressive engine (logic-bound), conservation kernel (bandwidth-bound)}
\end{axis}
\end{tikzpicture}
\caption{Modeled per-transaction cost vs.\ control-plane complexity
(analytical, from the cost model of \cref{sec:kernel}; not a measured
benchmark). An expressive engine dispatches on business cases on the hot
path, so its per-transaction cost grows with the rule set; the
conservation kernel's hot path is a branch-minimal integer fold over
fixed-width records, so its cost is a constant set by the record width.
The gap is structural, not a tuning artifact.}
\label{fig:throughput}
\end{figure}
